\subsection{Essential valuations of a Krull domain}

Let A be a Krull domain and K its field of fractions. The valuations defined by formula (4) of no. 3 (for \(x \in \mathbf{K}^*\) ) are called the essential valuations of K (or A).

We have remarked in the course of the proof of Theorem 2 that the valuations \(v_{\mathbb{P}}\) satisfy properties \((\mathrm{AK}_{\mathbb{I}}), (\mathrm{AK}_{\mathbb{II}})\) and \((\mathrm{AK}_{\mathbb{III}})\) of Definition 3. Moreover, these discrete valuations \(v_{\mathbb{P}}\) are normed: for every extremal divisor \(\mathbf{P} \in \mathbf{P}(\mathbf{A})\) , \(\mathbf{P} < 2\mathbf{P}\) and hence, if \(\mathfrak{a}\) and \(\mathfrak{b}\) are the divisorial ideals corresponding to \(\mathbf{P}\) and \(2\mathbf{P}\) , then \(\mathfrak{a} \supset \mathfrak{b}\) and \(\mathfrak{a} \neq \mathfrak{b}\) ; for \(x \in \mathfrak{a} - \mathfrak{b}\) , \(\operatorname{div}(x) \geqslant \mathbf{P}\) and \(\operatorname{div}(x) \not\geqslant 2\mathbf{P}\) , whence \(v_{\mathbb{P}}(x) = 1\) , which proves our assertion.

PROPOSITION 5. Let A be a Krull domain, K its field of fractions and \((v_{\mathbf{P}})_{\mathbf{P} \in \mathbf{P}(\mathbf{A})}\) the family of its essential valuations. Let \((n_{\mathbf{P}})_{\mathbf{P} \in \mathbf{P}(\mathbf{A})}\) be a family of rational integers which are zero except for a finite number of indices. Then the set of \(x \in \mathbf{K}\) such that \(v_{\mathbf{P}}(x) \geqslant n_{\mathbf{P}}\) for all \(\mathbf{P} \in \mathbf{P}(\mathbf{A})\) is the divisorial ideal \(\mathfrak{a}\) of \(\mathbf{A}\) such that \(\operatorname{div} \mathfrak{a} = \sum_{\mathbf{P} \in \mathbf{P}(\mathbf{A})} n_{\mathbf{P}} \cdot \mathbf{P}\) .

Let \(x \in K^*\) . In order that \(x \in a\) , it is necessary and sufficient that \(Ax \subset a\) , hence that \(\operatorname{div}(x) \geqslant \operatorname{div} a\) and hence, by (4), that \(v_{\mathbf{P}}(x) \geqslant n_{\mathbf{P}}\) for all \(\mathbf{P} \in \mathbf{P}(\mathbf{A})\) .

PROPOSITION 6. Let A be a Krull domain, K its field of fractions, \((v_{\iota})_{\iota \in I}\) a family of valuations on K with the properties of Definition 3 and \(\mathbf{A}_{\iota}\) the ring of \(v_{\iota}\) . Let S be a multiplicative subset of A not containing 0 and J the set of indices \(\iota \in I\) such that \(v_{\iota}\) is zero on S. Then \(\mathrm{S}^{-1}\mathrm{A} = \bigcap_{\iota \in I}\mathrm{A}_{\iota}\) ; in particular \(\mathrm{S}^{-1}\mathrm{A}\) is a Krull domain.

We write \(\mathbf{B} = \bigcap_{\iota \in \mathbf{J}} \mathbf{A}_{\iota}\) . Then \(\mathbf{S}^{-1} \subset \mathbf{B}\) and \(\mathbf{A} \subset \mathbf{B}\) and hence \(\mathbf{S}^{-1}\mathbf{A} \subset \mathbf{B}\) . Conversely, let \(x \in \mathbf{B}\) . Let \(\mathbf{J}'\) denote the finite set of indices \(\iota\) such that \(v_{\iota}(x) < 0\) . If \(\iota \in \mathbf{J}'\) , then \(x \notin \mathbf{A}_{\iota}\) , hence \(\iota \notin \mathbf{J}\) and hence there exists \(s_{\iota} \in \mathbf{S}\) such that \(v_{\iota}(s_{\iota}) > 0\) . Let \(n(\iota)\) be an integer \(>0\) such that \(v_{\iota}(s_{\iota}^{n(\iota)}x) \geqslant 0\) ; we write \(s = \prod_{\iota \in \mathbf{J}} s_{\iota}^{n(\iota)}\) . Then \(v_{\iota}(sx) \geqslant 0\) for all \(\iota \in \mathbf{I}\) and hence \(sx \in \mathbf{A}\) and \(x \in \mathbf{S}^{-1}\mathbf{A}\) . Thus \(\mathbf{B} = \mathbf{S}^{-1}\mathbf{A}\) .

COROLLARY 1. Let P be an extremal divisor of A and p the corresponding divisorial ideal. Then p is prime, the ring of \(v_{\mathbf{P}}\) is \(A_{\mathfrak{p}}\) and the residue field of \(v_{\mathbf{P}}\) is identified with the field of fractions of A/p.

Let \(S = A - p\) . By Proposition 5, \(v_{P}\) is zero on \(S\) and \(>0\) on \(p\) . Hence \(p\) is the intersection of \(A\) and the ideal of \(v_{P}\) and is therefore prime. On the other hand, for every extremal divisor \(Q \neq P\) , \(Q \not\geq P\) and hence the divisorial ideal \(q\) corresponding to \(Q\) is not contained in \(p\) ; thus \(q \cap S \neq \varnothing\) and hence, by Proposition 5, \(v_{Q}\) is not zero on \(S\) . This being so, the corollary follows from Proposition 6 and Chapter II, § 3, no. 1, Proposition 3.

COROLLARY 2. Let A be a Krull domain, K its field of fractions and \((v_{\iota})_{\iota\in I}\) a family of valuations with the properties of Definition 3. Then every essential valuation of A is equivalent to one of the \(v_{\iota}\) .

Let P be an extremal divisor of A and p the corresponding divisorial ideal. By Corollary 1, Proposition 5, Lemma 1 and assertion (1) in the proof of Theorem 2, no. 3, there exists \(\iota \in I\) such that the ring \(A_{\iota}\) of \(v_{\iota}\) contains the ring \(A_{p}\) of \(v_{P}\) . As \(v_{\iota}\) and \(v_{P}\) are of height 1, they are therefore equivalent (Chapter VI, §4, no. 5, Proposition 6).

PROPOSITION 7. Let A be a Krull domain, \((v_{\mathbb{P}})_{\mathbb{P} \in \mathbb{P}(\mathbb{A})}\) the family of its essential valuations and \(\mathfrak{a} \in \mathbf{I}(\mathbf{A})\) . Then the coefficient of P in div \(\mathfrak{a}\) is \(\inf_{y \in \mathfrak{a}} (v_{\mathbb{P}}(y))\) . If \(\mathfrak{p}\) is the divisorial prime ideal corresponding to the extremal divisor P, then \(\mathfrak{a}\mathrm{A}_{\mathfrak{p}} = \tilde{\mathfrak{a}}\mathrm{A}_{\mathfrak{p}}\) .

As \(\mathfrak{a} = \sum_{y \in \mathfrak{a}} Ay\) , Proposition 2 (b) (no. 1) shows that \(\text{div}(\mathfrak{a}) = \inf_{y \in \mathfrak{a}} (\text{div}(Ay))\) , whence our first assertion. The second follows immediately, since \(\text{div} \tilde{\mathfrak{a}} = \text{div} \mathfrak{a}\) and \(A_{\mathfrak{p}}\) is the ring of the discrete valuation \(v_{\mathfrak{p}}\) .

PROPOSITION 8. Let A be an integrally closed Noetherian domain.

\begin{enumerate}
\def\labelenumi{(\alph{enumi})}
\tightlist
\item
  Let \(\mathbf{P}\) be an extremal divisor of \(\mathbf{A}\) and \(\mathfrak{p}\) the corresponding divisorial prime ideal;
\end{enumerate}

for \(n \in \mathbf{N}\) , let \(\mathfrak{p}^{(n)} = \mathfrak{p}^n\mathbf{A}_{\mathfrak{p}} \cap \mathbf{A}\) ; then \(\mathfrak{p}^{(n)}\) is the set of \(x \in \mathbf{A}\) such that \(v_{\mathfrak{p}}(x) \geqslant n\) and is a \(\mathfrak{p}\)-primary ideal.

\begin{enumerate}
\def\labelenumi{(\alph{enumi})}
\setcounter{enumi}{1}
\tightlist
\item
  Let \(\mathfrak{a}\) be a divisorial integral ideal, \(n_1\mathrm{P}_1 + \dots + n_r\mathrm{P}_r\) the divisor of \(\mathfrak{a}\) (the \(\mathbf{P}_i\) being distinct extremal divisors) and \(\mathfrak{p}_i\) the divisorial prime ideal corresponding to \(\mathbf{P}_i\) .
\end{enumerate}

Then \(\mathfrak{a} = \bigcap_{i=1}^{r} \mathfrak{p}_i^{(n_i)}\) is the unique reduced primary decomposition of \(\mathfrak{a}\) and the \(\mathfrak{p}_i\) are not immersed.

By Corollary 1 to Proposition 6, the relation \(x \in \mathfrak{p}^n\mathrm{A}_{\mathfrak{p}} = (\mathfrak{p}\mathrm{A}_{\mathfrak{p}})^n\) is equivalent to \(v_{\mathfrak{p}}(x) \geqslant n\) ; on the other hand, as \(\mathrm{A}_{\mathfrak{p}}\) is a discrete valuation ring, \((\mathfrak{p}\mathrm{A}_{\mathfrak{p}})^n\) is \((\mathfrak{p}\mathrm{A}_{\mathfrak{p}})\) -primary (Chapter IV, § 2, no. 1, Example 4) and hence \(\mathfrak{p}^{(n)}\) is \(\mathfrak{p}\) -primary (Chapter IV, § 2, no. 1, Proposition 3); this shows (a). Proposition 5 certainly shows that \(\mathfrak{a} = \bigcap_{i=1}^{r} \mathfrak{p}_i^{(n_i)}\) . As \(\mathfrak{p}_i \not\subset \mathfrak{p}_j\) for \(i \neq j\) this primary decomposition is reduced: For if \(\mathfrak{p}_i^{(n_i)} \supset \bigcap_{j \neq i} \mathfrak{p}_j^{(n_j)} \supset \prod_{j \neq i} \mathfrak{p}_j^{(n_j)}, \mathfrak{p}_i\) would contain one of the \(\mathfrak{p}_j\) for \(j \neq i\) (Chapter II, § 1, no. 1, Proposition 1). The uniqueness follows from Chapter IV, § 2, no. 3, Proposition 5.

